% Diagrama F — trajetória da relação efetiva de vencedor W_ij ao longo de n:
% empate inicial não resolvido, primeiro vencedor, empate posterior que
% mantém o vencedor, e inversão válida que atualiza W e define uma célula em R.
\documentclass[tikz,border=6pt]{standalone}
% Preâmbulo compartilhado dos diagramas conceituais do relatório CoInfoSim.
% Compilação: xelatex (determinística; nenhum dado experimental é usado).
\usepackage{fontspec}
\setmainfont{Liberation Sans}
\setsansfont{Liberation Sans}
\usepackage{amsmath}
\usepackage{bm}
\usepackage{tikz}
\usetikzlibrary{arrows.meta,positioning,calc,fit,backgrounds,decorations.pathreplacing,matrix,patterns}
\definecolor{CoinAccent}{RGB}{25,80,140}
\definecolor{CoinAccentDark}{RGB}{18,55,95}
\definecolor{CoinAccentPale}{RGB}{237,242,248}
\definecolor{CoinGray}{RGB}{95,95,95}
\definecolor{CoinGrayPale}{RGB}{246,246,246}
\definecolor{CoinWin}{RGB}{76,140,90}    % verde discreto (vitória)
\definecolor{CoinLose}{RGB}{182,88,79}   % vermelho discreto (derrota)
\tikzset{
  bloco/.style={draw=CoinAccentDark, fill=CoinAccentPale, rounded corners=1.2pt,
                align=center, font=\small, inner sep=4pt, minimum height=7mm},
  bloconeutro/.style={draw=CoinGray, fill=CoinGrayPale, rounded corners=1.2pt,
                align=center, font=\small, inner sep=4pt, minimum height=7mm},
  blocoreal/.style={draw=CoinAccentDark, fill=white, rounded corners=1.2pt,
                align=center, font=\small, inner sep=4pt, minimum height=7mm},
  seta/.style={-{Stealth[length=2.4mm]}, thick, CoinAccentDark},
  setacinza/.style={-{Stealth[length=2.4mm]}, thick, CoinGray},
  rotulo/.style={font=\footnotesize\itshape, text=CoinGray, align=center},
}

\begin{document}
\begin{tikzpicture}[node distance=6mm]
  \def\ngrid{2/{$n_1{=}2$}, 1/{$n_2{=}4$}, 0/{$n_3{=}8$}, -1/{$n_4{=}16$}, -2/{$n_5{=}32$}}
  \def\xstep{2.35}

  % Eixo horizontal (grade de n)
  \draw[->, CoinGray] (-0.6,0) -- (11.0,0) node[below right, rotulo] {$n$ (grade de potências de 2)};
  \foreach \i/\lab [count=\k from 0] in {2/{$n_1{=}2$},1/{$n_2{=}4$},0/{$n_3{=}8$},-1/{$n_4{=}16$},-2/{$n_5{=}32$}}{
    \pgfmathsetmacro{\x}{\k*\xstep}
    \draw[CoinGray] (\x,0.08) -- (\x,-0.08);
    \node[rotulo] at (\x,-0.42) {\lab};
  }

  % Estados de U_ij e W_ij em cada ponto (0=empate exato / observado; +1 = S_i vence; -1 = S_j vence)
  % n1: empate exato inicial -> W indefinido (0)
  % n2: S_i vence estritamente -> primeiro vencedor, inicializa W (NÃO é inversão)
  % n3: empate exato -> W mantém S_i (carry-forward; NÃO é inversão)
  % n4: S_j vence estritamente -> INVERSÃO VÁLIDA (W_{n3}=S_i != 0, W_{n4}=S_j != 0, muda de sinal)
  % n5: S_j continua vencendo -> W mantém S_j
  \foreach \k/\y/\txt/\cor in {
    0/0/{$U{=}0$\\vazio, indefinido}/CoinGray,
    1/1.35/{$U{=}{+1}$\\primeiro vencedor}/CoinWin,
    2/1.35/{$U{=}0$\\empate: mantém $S_i$}/CoinAccent,
    3/-1.35/{$U{=}{-1}$\\\textbf{inversão válida}}/CoinLose,
    4/-1.35/{$U{=}{-1}$\\mantém $S_j$}/CoinLose
  }{
    \pgfmathsetmacro{\x}{\k*\xstep}
    \node[draw=\cor, fill=\cor!15, rounded corners=1pt, font=\scriptsize, align=center,
          text width=2.05cm, inner sep=2pt, minimum height=9mm] (pt\k) at (\x,\y) {\txt};
  }

  % Linha de base do vencedor efetivo W_ij (S_i ou S_j) sob os pontos
  \node[rotulo, font=\tiny] at (0,-1.0) {$W_{ij}$: indefinido};
  \node[rotulo, font=\tiny] at (\xstep,-1.0) {$W_{ij}=S_i$};
  \node[rotulo, font=\tiny] at (2*\xstep,-1.0) {$W_{ij}=S_i$};
  \node[rotulo, font=\tiny] at (3*\xstep,-1.0) {$W_{ij}=S_j$};
  \node[rotulo, font=\tiny] at (4*\xstep,-1.0) {$W_{ij}=S_j$};

  % Conectores tracejados destacando a transição de inversão entre n3 e n4
  \draw[CoinGray, dashed] (pt1.south) -- (pt2.north);
  \draw[CoinLose, thick] (pt2.south) -- (pt3.north)
    node[midway, right=1mm, rotulo, text=CoinLose, font=\tiny] {inversão};
  \draw[CoinGray, dashed] (pt3.south) -- (pt4.north);

  % Célula resultante em R
  \node[draw=CoinAccentDark, thick, rounded corners=1pt, fill=CoinAccentPale, font=\small,
        align=center, text width=5.4cm, inner sep=5pt] (rcell) at (5.0,-3.1)
    {$R_{ij}(n_5)=16$\\ \normalfont\footnotesize única célula de $\mathbf{R}$ para o par $(i,j)$: o tamanho amostral da última inversão válida observada};
  \draw[seta] (pt3.south) -- ++(0,-0.4) -| (rcell.north);

  \node[rotulo, text width=9.5cm] at (5.0,-4.1)
    {$n_1$: empate exato inicial, indefinido. $\quad n_2$: primeiro vencedor estrito, inicializa $W_{ij}$ (não é inversão). $\quad n_3$: empate exato posterior mantém o vencedor anterior (não é inversão). $\quad n_4$: vencedor definido substituído pelo outro --- inversão válida. $\quad n_5$: vencedor mantido.};
\end{tikzpicture}
\end{document}
